\section{Suites \texorpdfstring{$M$}{M}-régulières et suites
\texorpdfstring{$\mathscr F$}{F}-régulières}
\label{section:0.15-fr}

\subsection{Suites \texorpdfstring{$M$}{M}-régulières et suites
\texorpdfstring{$M$}{M}-quasi-régulières}
\label{subsection:0.15.1-fr}
\label{0.15.1-fr}

\begin{env}[15.1.1]
\label{0.15.1.1-fr}
\oldpage[0\textsubscript{IV-1}]{12}
Soient $A$ un anneau, $N$ un $A$-module~; dans ce qui suit, nous noterons
$N[T_1,\ldots,T_n]$, pour abréger, le $A$-module gradué
$N\otimes_A A[T_1,\ldots,T_n]$ ($T_i$ indéterminées), où $N$ est gradué par la
graduation triviale et $A[T_1,\ldots,T_n]$ par le degré total.

Soient $M$ un $A$-module, $f_i$ ($1\leq i\leq n$) des éléments de $A$,
$\mathfrak{J}=\sum_{i=1}^n f_iA$ l'idéal de $A$ engendré par les $f_i$, et
munissons $M$ de la filtration $\mathfrak{J}$-préadique~; on définit alors un
\emph{homomorphisme gradué surjectif de $A$-modules gradués}, de degré $0$,
\[
  \varphi : (\gr_0(M))[T_1,\ldots,T_n]\longrightarrow\gr_\bullet(M)
  \tag{15.1.1.1}\label{0.15.1.1.1-fr}
\]
de la façon suivante~: si $\xi\in\gr_0(M)$ est la classe mod.
$\mathfrak{J}M$ d'un élément $x\in M$, au couple formé de $\xi$ et d'un
polynôme $P[T_1,\ldots,T_n]$ homogène de degré $k$, on fait correspondre la
classe mod. $\mathfrak{J}^{k+1}M$ de $P(f_1,\ldots,f_n)x$~; il est immédiat
que cette classe
\oldpage[0\textsubscript{IV-1}]{13}
ne dépend que de $\xi$ et de $P$, et si $S_k$ est l'ensemble des polynômes
homogènes de degré $k$ dans $A[T_1,\ldots,T_n]$, on a ainsi défini une
application $A$-linéaire
\[
  \varphi_k:\gr_0(M)\otimes_A S_k\longrightarrow\gr_k(M),
\]
qui est surjective par définition de $\mathfrak{J}^kM$.

Nous dirons dans ce qui suit qu'un élément $f\in A$ est \emph{non diviseur de
zéro dans un $A$-module $M$} si l'homothétie $f_M:x\mapsto fx$ de $M$ est
\emph{injective}.
\end{env}

\begin{lemma}[15.1.2]
\label{0.15.1.2-fr}
Soient $M$ un $A$-module, $f$ un élément de $A$, et munissons $M$ de la
filtration $(fA)$-préadique. Si $f$ est non diviseur de zéro dans $M$,
l'homomorphisme canonique défini dans (15.1.1)~:
\[
  \varphi:(M/fM)[T]\longrightarrow\gr_\bullet(M)
  \tag{15.1.2.1}\label{0.15.1.2.1-fr}
\]
est bijectif. La réciproque est vraie si $M$ est séparé pour la topologie
$(fA)$-préadique.
\end{lemma}

\begin{proof}
On a ici $\gr_k(M)=f^kM/f^{k+1}M$~; dire que $\varphi$ est injective signifie
que pour tout $x\in M$ et tout $k$, la relation $f^kx\in f^{k+1}M$ implique
$x\in fM$~; il est clair qu'il en est bien ainsi si $f$ est non diviseur de
zéro dans $M$. Inversement, remarquons que l'homothétie
$f_M:x\mapsto fx$ définit par passage aux quotients, des homomorphismes
$\gr_k(M)\to\gr_{k+1}(M)$, donc un homomorphisme gradué $g$ de degré $1$ de
$\gr_\bullet(M)$ dans lui-même~; et dire que $\varphi$ est injectif exprime
encore que $g$ est injectif. Si on suppose la filtration de $M$ séparée, il en
résulte que $f_M$ est injective (Bourbaki, \emph{Alg.~comm.}, chap.~III,
\S~2, no~8, cor.~1 du th.~1).
\end{proof}

\begin{corollary}[15.1.3]
\label{0.15.1.3-fr}
Soient $A$ un anneau noethérien, $M$ un $A$-module de type fini. Pour que $f$
soit non diviseur de zéro dans $M$, il faut et il suffit que $\varphi$ soit
bijectif.
\end{corollary}

\begin{proof}
On doit seulement démontrer que la condition est suffisante~; pour prouver
que $f_M$ est injective, il suffit de montrer que pour tout idéal premier
$\mathfrak p$ de $A$, l'homothétie de rapport
$f'=f/1\in A'=A_{\mathfrak p}$ dans le $A_{\mathfrak p}$-module
$M_{\mathfrak p}=M'$ est injective (Bourbaki, \emph{Alg.~comm.}, chap.~II,
\S~3, no~4, th.~1). Or on a
$f'^kM'/f'^{k+1}M'=(f^kM/f^{k+1}M)\otimes_A A'$, et il est immédiat que
l'homomorphisme canonique
$\varphi':(M'/f'M')[T]\to\gr_\bullet(M')$ est égal à
$\varphi\otimes 1_{A'}$~; si $\varphi$ est injectif, il en est donc de même de
$\varphi'$ ($0_{\mathrm I}$, 1.3.2). On est donc ramené au cas où $A$ est un
anneau noethérien \emph{local}, dont nous désignerons par $\mathfrak m$
l'idéal maximal. Si $f\notin\mathfrak m$, alors $f$ est inversible et il n'y a
rien à démontrer. Si $f\in\mathfrak m$, on sait que $M$ est séparé pour la
topologie $(fA)$-préadique ($0_{\mathrm I}$, 7.3.5), et il suffit donc
d'appliquer (15.1.2).
\end{proof}

\begin{definition}[15.1.4]
\label{0.15.1.4-fr}
Soient $A$ un anneau, $M$ un $A$-module. On dit qu'un élément $f\in A$ est
\emph{$M$-régulier} (resp. \emph{$M$-quasi-régulier}) si $f$ est non diviseur
de zéro dans $M$ (resp. si l'homomorphisme (15.1.2.1) est bijectif).
\end{definition}

Nous venons donc de démontrer que tout élément $M$-régulier est
$M$-quasi-régulier, et que la réciproque est vraie si $A$ est noethérien et
$M$ un $A$-module de type fini. Lorsqu'on ne suppose plus $M$ de type fini,
cette réciproque peut être en défaut~: par exemple, si $A$ est un anneau
principal qui n'est pas un corps, $K$ son corps des fractions et $M$ le
$A$-module $K/A$, on a $fM=M$ pour tout $f\neq0$ dans $A$, donc les deux
membres de (15.1.2.1) sont réduits à $0$, mais $f$ est alors diviseur de zéro
dans $M$.

\begin{env}[15.1.5]
\label{0.15.1.5-fr}
\oldpage[0\textsubscript{IV-1}]{14}
Supposons maintenant $M$ muni d'une \emph{filtration} $(M_k)$ telle que
$M_0=M$, et désignons par
$\gr_\bullet(M)=\bigoplus_{k\geq0}M_k/M_{k+1}$ le $A$-module gradué associé.
Les notations étant celles de (15.1.1), posons, pour tout $k\geq0$,
\[
  M'_k=M_k+\mathfrak{J}M_{k-1}+\cdots+
  \mathfrak{J}^{k-1}M_1+\mathfrak{J}^kM_0.
  \tag{15.1.5.1}\label{0.15.1.5.1-fr}
\]
Il est clair que $(M'_k)$ est une filtration sur $M$. Lorsque
$M_k=\mathfrak{Q}^kM$, où $\mathfrak Q$ est un idéal de $A$, la filtration
$(M'_k)$ n'est autre que la filtration
$(\mathfrak{J}+\mathfrak{Q})$-préadique. Nous noterons
$\gr'_\bullet(M)$ le $A$-module gradué associé à $(M'_k)$. Notons encore
\[
  (\gr_\bullet(M)\otimes_A(A/\mathfrak{J}))[T_1,\ldots,T_n]
\]
le produit tensoriel de $A$-modules gradués ($\mathrm{II}$, 2.1.2)
\[
  (\gr_\bullet(M)\otimes_A(A/\mathfrak{J}))
  \otimes_A A[T_1,\ldots,T_n].
\]
On définit un \emph{homomorphisme gradué surjectif de $A$-modules gradués}, de
degré $0$,
\[
  \psi:(\gr_\bullet(M)\otimes_A(A/\mathfrak{J}))[T_1,\ldots,T_n]
  \longrightarrow\gr'_\bullet(M)
  \tag{15.1.5.2}\label{0.15.1.5.2-fr}
\]
de la façon suivante~: si $\alpha\in A/\mathfrak{J}$ est la classe d'un
élément $a\in A$, $\xi\in M_k/M_{k+1}$ la classe d'un élément $x\in M_k$ et
$P[T_1,\ldots,T_n]$ un polynôme homogène de degré $h$, on fait correspondre au
triplet $(\alpha,\xi,P)$ la classe mod. $M'_{h+k+1}$ de
$P(f_1,\ldots,f_n)ax$~; on vérifie encore aussitôt que cet élément ne dépend
bien que de $\xi$, $\alpha$ et $P$, et cela définit l'homomorphisme gradué
(15.1.5.2), qui est surjectif en vertu de la définition (15.1.5.1). On
retrouve l'homomorphisme (15.1.1.1) en considérant la filtration sur $M$ telle
que $M_1=0$.
\end{env}

\begin{lemma}[15.1.6]
\label{0.15.1.6-fr}
On garde les notations de (15.1.5) et on suppose $n=1$, en posant $f=f_1$. Si
$f$ est $\gr_\bullet(M)$-régulier, l'homomorphisme (15.1.5.2) est bijectif.
\end{lemma}

\begin{proof}
Soit $Q_k$ (resp. $Q'_k$) le sous-$A$-module des termes de degré $k$ au
premier (resp. second) membre de (15.1.5.2). Munissons $Q_k$ de la filtration
\[
  (Q_k)_i=\sum_{j\leq k-i}
  (\gr_{k-j}(M)\otimes_A(A/fA))T^j
  \tag{15.1.6.1}\label{0.15.1.6.1-fr}
\]
de sorte que $(Q_k)_0=Q_k$ et $(Q_k)_{k+1}=0$, et pour cette filtration, on a
\[
  \gr_i(Q_k)=(\gr_i(M)\otimes_A(A/fA))T^{k-i}.
\]
Munissons $Q'_k$ de la filtration formée des
$\psi((Q_k)_i)=(Q'_k)_i$. Comme ces filtrations sont \emph{finies}, il suffit
de prouver que les homomorphismes
$\psi_{ki}:\gr_i(Q_k)\to\gr_i(Q'_k)$ déduits de $\psi_k$ sont injectifs
(Bourbaki, \emph{Alg.~comm.}, chap.~III, \S~3, no~8, cor.~1 du th.~1). Or on
a
\[
  \gr_i(Q_k)=((M_i/M_{i+1})\otimes_A(A/fA))T^{k-i}~;
\]
d'autre part, $(Q'_k)_{i+1}$ est l'image de
\[
  M_k+fM_{k-1}+\cdots+f^{k-i-1}M_{i+1}
\]
\oldpage[0\textsubscript{IV-1}]{15}
dans $M'_k/M'_{k+1}$. On est donc ramené à écrire que, pour $x\in M_i$, la
relation
\[
  f^{k-i}x\in M_k+fM_{k-1}+\cdots+
  f^{k-i-1}M_{i+1}+M'_{k+1}
  \tag{15.1.6.2}\label{0.15.1.6.2-fr}
\]
implique $x\in fM_i+M_{i+1}$. Or, le second membre de (15.1.6.2) est contenu
dans $M_{i+1}+M'_{k+1}$, et en vertu de (15.1.5.1), $M'_{k+1}$ est contenu
dans $M_{i+1}+f^{k+1-i}M$. L'hypothèse sur $f$ entraîne que $f$ est non
diviseur de zéro dans $M/M_h$ pour tout $h$ (Bourbaki, \emph{loc.~cit.}). La
relation $f^{k-i}x\in f^{k+1-i}M+M_{i+1}$ entraîne donc (en raisonnant dans
$M/M_{i+1}$) l'existence d'un $y\in M$ tel que $x-fy\in M_{i+1}$, et comme
$x\in M_i$, on a $fy\in M_i$, d'où (en raisonnant dans $M/M_i$), $y\in M_i$,
et finalement $x\in fM_i+M_{i+1}$. C.Q.F.D.
\end{proof}

\begin{definition}[15.1.7]
\label{0.15.1.7-fr}
Soit $(f_i)_{1\leq i\leq n}$ une suite d'éléments de $A$. On dit qu'elle est
\emph{$M$-régulière} si, pour $1\leq i\leq n$, $f_i$ est
$\left(M/(\sum_{1\leq j\leq i-1}f_jM)\right)$-régulier. On dit que la suite
$(f_i)$ est \emph{$M$-quasi-régulière} si
l'homomorphisme canonique $\varphi$ défini dans (15.1.1.1) est bijectif.

Lorsque $M=A$, on dit simplement «~suite régulière~» et
«~suite quasi-régulière~».
\end{definition}

\begin{proposition}[15.1.8]
\label{0.15.1.8-fr}
Les notations étant celles de (15.1.5), on suppose que la suite
$(f_i)_{1\leq i\leq n}$ est $\gr_\bullet(M)$-régulière. Alors
l'homomorphisme canonique (15.1.5.2) est bijectif.
\end{proposition}

\begin{proof}
La proposition n'est autre que (15.1.6) pour $n=1$~; raisonnons par récurrence
sur $n$. Posons
\[
  \mathfrak{J}''=\sum_{i=1}^{n-1}f_iA,
\]
de sorte que $\mathfrak{J}=\mathfrak{J}''+f_nA$, et soit
\[
  M''_k=M_k+\mathfrak{J}''M_{k-1}+\cdots+
  \mathfrak{J}''^{k-1}M_1+\mathfrak{J}''^kM_0~;
\]
on peut donc écrire
\[
  M'_k=M''_k+f_nM''_{k-1}+\cdots+f_n^{k-1}M''_1+f_n^kM''_0.
\]
Notons $\gr''_\bullet(M)$ le $A$-module gradué associé à la filtration
$(M''_k)$ de $M$. On peut écrire, à des isomorphismes canoniques près
\[
  (\gr_\bullet(M)\otimes(A/\mathfrak{J}))\otimes_A A[T_1,\ldots,T_n]
  =
  (\gr_\bullet(M)\otimes(A/\mathfrak{J}''))[T_1,\ldots,T_{n-1}]
  \otimes(A/f_nA)\otimes_A A[T_n]
\]
et par suite (15.1.5.2) se factorise en
\[
  (\gr''_\bullet(M)\otimes(A/f_nA))[T_n]\longrightarrow\gr'_\bullet(M)
  \tag{15.1.8.1}\label{0.15.1.8.1-fr}
\]
et
\[
  ((\gr_\bullet(M)\otimes(A/\mathfrak{J}''))[T_1,\ldots,T_{n-1}])
  \otimes(A/f_nA)\otimes_A A[T_n]
  \xrightarrow{\psi'\otimes1}
  \gr''_\bullet(M)\otimes(A/f_nA)\otimes_A A[T_n]
\]
$\psi'$ étant l'application (15.1.5.2) où l'on remplace $n$ par $n-1$,
$\mathfrak{J}$ par $\mathfrak{J}''$ et $\gr'_\bullet(M)$ par
$\gr''_\bullet(M)$. L'hypothèse de récurrence entraîne que $\psi'$ est
bijective et il reste donc à voir qu'il en est de même de (15.1.8.1). En
appliquant (15.1.6), il suffit de montrer que $f_n$ est
$\gr''_\bullet(M)$-régulier. Mais $\gr''_\bullet(M)$ s'identifie à
\[
  (\gr_\bullet(M)\otimes(A/\mathfrak{J}''))
  [T_1,\ldots,T_{n-1}]~;
\]
par hypothèse, l'homothétie de rapport $f_n$ dans
$\gr_\bullet(M)\otimes(A/\mathfrak{J}'')$ est bijective~; il en est par suite
de même de l'homothétie de rapport $f_n$ dans $\gr''_\bullet(M)$, puisque
$A[T_1,\ldots,T_{n-1}]$ est un $A$-module libre.
\end{proof}

\begin{proposition}[15.1.9]
\label{0.15.1.9-fr}
Soient $M$ un $A$-module, $(f_i)_{1\leq i\leq n}$ une suite d'éléments de $A$.
Si la suite $(f_i)$ est $M$-régulière, elle est $M$-quasi-régulière. La
réciproque est vraie lorsqu'on suppose en outre que les $A$-modules
\[
  M,\quad M/f_1M,\quad\ldots,\quad
  M/\left(\sum_{1\leq j\leq n-1}f_jM\right)
\]
sont séparés pour la topologie $\mathfrak{J}$-préadique (où
$\mathfrak{J}=\sum_{i=1}^n f_iA$).
\end{proposition}

\begin{proof}
La première assertion résulte de (15.1.8), où l'on prend $M_1=0$, et par suite
l'homomorphisme (15.1.5.2) se réduit à (15.1.1.1). Inversement, supposons que
l'homomorphisme $\varphi$ de (15.1.1.1) soit bijectif. Notons que si on pose
\[
  \mathfrak{J}''=\sum_{i=1}^{n-1}f_iA,
\]
\oldpage[0\textsubscript{IV-1}]{16}
et si on désigne par $\gr''_\bullet(M)$ le $A$-module gradué associé à $M$
filtré par la filtration $\mathfrak{J}''$-préadique, $\varphi$ se factorise en
\[
  (\gr''_\bullet(M)\otimes(A/f_nA))[T_n]\longrightarrow\gr_\bullet(M)
  \tag{15.1.9.1}\label{0.15.1.9.1-fr}
\]
et en $\varphi'\otimes1$, où $\varphi'$ est l'homomorphisme canonique
\[
  (\gr''_0(M))[T_1,\ldots,T_{n-1}]
  \longrightarrow\gr''_\bullet(M).
\]
Comme ces deux homomorphismes sont surjectifs, ils sont bijectifs si
$\varphi$ est supposé bijectif, donc $\varphi'$ est nécessairement injectif,
et par suite bijectif puisqu'il est surjectif. On peut alors raisonner par
récurrence sur $n$ (le cas $n=1$ résultant de (15.1.2)), car sur
\[
  M,\quad M/f_1M,\quad\ldots,\quad
  M/\left(\sum_{j=1}^{n-2}f_jM\right)
\]
la topologie $\mathfrak{J}''$-préadique est plus fine que la topologie
$\mathfrak{J}$-préadique, et par suite est séparée. On voit donc que la suite
$(f_i)_{1\leq i\leq n-1}$ est $M$-régulière. D'autre part, montrons que
l'hypothèse entraîne que si $y\in M/\mathfrak{J}''M$ est tel que
\[
  f_n^ky\in f_n^{k+1}(M/\mathfrak{J}''M),
\]
alors on a $y\in f_n(M/\mathfrak{J}''M)$. En effet, si $x$ est un représentant
de $y$, montrons que l'on a $f_n^kx\in\mathfrak{J}^{k+1}M$~; notre assertion
résultera alors du fait que (15.1.9.1) est injectif. L'hypothèse entraîne
\[
  f_n^kx\in f_n^{k+1}M+\mathfrak{J}''M\subset\mathfrak{J}M,
\]
d'où, puisque $\varphi$ est injectif,
$f_n^{k-1}x\in\mathfrak{J}^kM$~; par récurrence, on obtient ainsi
$x\in\mathfrak{J}M$, d'où finalement
$f_n^kx\in\mathfrak{J}^{k+1}M$. On termine le raisonnement comme dans
(15.1.2)~: sur $M/\mathfrak{J}''M$, la filtration
$\mathfrak{J}$-préadique est identique à la filtration
$(f_nA)$-préadique, donc est séparée par hypothèse, et on conclut que
l'homothétie de rapport $f_n$ dans $M/\mathfrak{J}''M$ est injective, ce qui
achève de montrer que la suite $(f_i)_{1\leq i\leq n}$ est $M$-régulière.
\end{proof}

\begin{corollary}[15.1.10]
\label{0.15.1.10-fr}
Avec les notations de (15.1.9), supposons que pour tout $A$-module $N$ de type
fini, $M\otimes_A N$ soit séparé pour la topologie
$\mathfrak{J}$-préadique. Alors, pour que la suite $(f_i)$ soit
$M$-régulière, il faut et il suffit qu'elle soit $M$-quasi-régulière.
\end{corollary}

On notera que l'hypothèse de (15.1.10) entraîne que, si la suite
$(f_i)_{1\leq i\leq n}$ est $M$-régulière, il en est de même de la suite
$(f_{\sigma(i)})$ pour toute permutation $\sigma$ de l'intervalle $[1,n]$.

\begin{corollary}[15.1.11]
\label{0.15.1.11-fr}
Soient $A$ un anneau noethérien, $M$ un $A$-module de type fini,
$(f_i)_{1\leq i\leq n}$ une suite finie d'éléments contenus dans le radical de
$A$. Alors, pour que la suite $(f_i)$ soit $M$-régulière, il faut et il suffit
qu'elle soit $M$-quasi-régulière.
\end{corollary}

\begin{proof}
C'est un cas particulier de (15.1.10), puisque tout $A$-module de type fini
est alors séparé pour la topologie $\mathfrak{J}$-préadique
($0_{\mathrm I}$, 7.3.5).
\end{proof}

\begin{remarks}[15.1.12]
\label{0.15.1.12-fr}
\mbox{}
\begin{enumerate}
  \item[(i)] Lorsque $A$ est un anneau local noethérien, $M$ un $A$-module de
  type fini, et lorsque les $f_i$ sont contenus dans l'idéal maximal de $A$,
  la notion de suite $M$-régulière a été introduite par J.-P.~Serre sous le
  nom de $M$-\emph{suite} [17].
  \item[(ii)] Même si $A$ est noethérien et $M=A$, une suite $(f,g)$ de deux
  éléments de $A$ qui est quasi-régulière n'est pas nécessairement régulière.
  Par exemple, la suite $(0,1)$ est quasi-régulière, puisqu'alors
  $\mathfrak{J}=fA+gA=A$, donc $\gr_\bullet(A)=0$~; mais si $A\neq0$ elle
  n'est pas régulière~; on notera toutefois que dans ce cas la suite $(1,0)$
  est régulière, car $1$ n'est pas diviseur de $0$ dans $A$, et comme
  $A/A=0$, $0$ n'est pas diviseur de zéro dans ce module (cf. (15.1.1)).
  \item[(iii)]
  \oldpage[0\textsubscript{IV-1}]{17}
  Dans (15.1.6), on ne peut en général remplacer l'hypothèse que $f$ est
  $\gr_\bullet(M)$-régulier par l'hypothèse qu'il est
  $\gr_\bullet(M)$-quasi-régulier. Pour le voir, reprenons l'exemple (15.1.4)
  où $A$ est un anneau principal de corps des fractions $K\neq A$ et prenons
  $M=K$, la filtration $(M_k)$ de $M$ étant définie par
  $M_0=M=K$, $M_1=A$, $M_k=0$ pour $k\geq2$, d'où
  $\gr_0(M)=K/A$, $\gr_1(M)=A$, $\gr_k(M)=0$ pour $k\geq2$. Si $f\neq0$ dans
  $A$, on a vu que $f$ n'est pas $\gr_\bullet(M)$-régulier, mais comme on a
  \[
    f^k\gr_\bullet(M)/f^{k+1}\gr_\bullet(M)=f^kA/f^{k+1}A,
  \]
  il est immédiat que l'homomorphisme (15.1.2.1) (où $M$ est remplacé par
  $\gr_\bullet(M)$) est injectif, donc $f$ est
  $\gr_\bullet(M)$-quasi-régulier. Mais avec les notations de (15.1.6), on a
  alors $M'_k=K$ pour tout $k\geq0$, donc $\gr'_\bullet(M)=0$, tandis que le
  premier membre de (15.1.5.2) n'est pas réduit à $0$.
\end{enumerate}
\end{remarks}

\begin{proposition}[15.1.13]
\label{0.15.1.13-fr}
Soient $A$ un anneau, $M$, $N$ deux $A$-modules, $(f_i)_{1\leq i\leq n}$ une
suite d'éléments de $A$. Si la suite $(f_i)$ est $M$-régulière et si $N$ est
un $A$-module plat, alors la suite $(f_i)$ est
$M\otimes_A N$-régulière. La réciproque est vraie si $N$ est un $A$-module
fidèlement plat.
\end{proposition}

\begin{proof}
En effet (sans hypothèse sur $N$),
\[
  (M\otimes_A N)/
  \left(\sum_{j=1}^{i-1}f_j(M\otimes_A N)\right)
\]
est isomorphe à
\[
  \left(M/\left(\sum_{j=1}^{i-1}f_jM\right)\right)\otimes_A N~;
\]
il suffit alors d'appliquer ($0_{\mathrm I}$, 6.1.1) et
($0_{\mathrm I}$, 6.4.1) à l'homothétie de rapport $f_i$ dans
$M/(\sum_{j=1}^{i-1}f_jM)$.
\end{proof}

\begin{corollary}[15.1.14]
\label{0.15.1.14-fr}
Soient $A$ un anneau, $M$ un $A$-module, $(f_i)_{1\leq i\leq n}$ une suite
finie d'éléments de $A$. Soient $\rho:A\to B$ un homomorphisme d'anneaux,
$M'=M_{(B)}=M\otimes_A B$, $f'_i=\rho(f_i)$ pour $1\leq i\leq n$. Si $B$ est
un $A$-module plat et si la suite $(f_i)$ est $M$-régulière, alors la suite
$(f'_i)$ est $M'$-régulière~; la réciproque est vraie si $B$ est un
$A$-module fidèlement plat.
\end{corollary}

\begin{proof}
En effet,
\[
  M'/\left(\sum_{j=1}^{i-1}f'_jM'\right)
\]
s'identifie à
\[
  \left(M/\left(\sum_{j=1}^{i-1}f_jM\right)\right)\otimes_A B,
\]
et l'homothétie de rapport $f'_i$ dans le $B$-module
$M'/(\sum_{j=1}^{i-1}f'_jM')$ à l'homothétie de rapport $f_i$ dans le
$A$-module
\[
  \left(M/\left(\sum_{j=1}^{i-1}f_jM\right)\right)\otimes_A B~;
\]
l'assertion résulte donc de (15.1.13).
\end{proof}

\begin{proposition}[15.1.15]
\label{0.15.1.15-fr}
Soient $A$ un anneau, $B$ une $A$-algèbre commutative,
$(f_i)_{1\leq i\leq n}$ une suite finie d'éléments de $B$, $M$ un
$B$-module. Supposons que la suite $(f_i)$ soit $M$-régulière et que chacun
des $A$-modules
\[
  M/\left(\sum_{j=1}^{i-1}f_jM\right)\quad(1\leq i\leq n)
\]
soit plat. Alors, pour tout homomorphisme d'anneaux $\rho:A\to A'$, si l'on
pose $B'=B\otimes_A A'$, $M'=M\otimes_A A'$ et $f'_i=f_i\otimes1$
($1\leq i\leq n$), la suite $(f'_i)$ est $M'$-régulière et les
$A'$-modules
\[
  M'/\left(\sum_{j=1}^{i-1}f'_jM'\right)\quad(1\leq i\leq n)
\]
sont plats.
\end{proposition}

\begin{proof}
Considérons la suite (exacte par hypothèse)
\[
  0\longrightarrow M\xrightarrow{f_1}M\longrightarrow M/f_1M
  \longrightarrow0
\]
(la flèche $M\xrightarrow{f_1}M$ étant l'homothétie de rapport $f_1$ dans
$M$). L'hypothèse que $M/f_1M$ est $A$-plat entraîne que la suite
\[
  0\longrightarrow M\otimes_A A'
  \xrightarrow{f_1\otimes1}M\otimes_A A'
  \longrightarrow(M/f_1M)\otimes_A A'\longrightarrow0
\]
\oldpage[0\textsubscript{IV-1}]{18}
est exacte ($0_{\mathrm I}$, 6.1.2). Cela montre que l'homothétie de rapport
$f'_1$ dans $M'$ est injective.

Posons ensuite, pour $1\leq i\leq n-1$,
\[
  M_i=M/\left(\sum_{j=1}^i f_jM\right),\qquad
  M'_i=M'/\left(\sum_{j=1}^i f'_jM'\right)~;
\]
on a $M'_i=M_i\otimes_A A'$, $M_i/f_{i+1}M_i=M_{i+1}$,
$M'_i/f'_{i+1}M'_i=M'_{i+1}$~; il suffit de remplacer dans le raisonnement
précédent $M$ et $f_1$ par $M_i$ et $f_{i+1}$ pour conclure que l'homothétie
de rapport $f'_{i+1}$ dans $M'_i$ est injective, en raison de la platitude de
$M_i/f_{i+1}M_i$. La dernière assertion résulte de ($0_{\mathrm I}$, 6.2.1).
\end{proof}

\begin{proposition}[15.1.16]
\label{0.15.1.16-fr}
Soient $A$, $B$ deux anneaux locaux noethériens, $k$ le corps résiduel de
$A$, $\varphi:A\to B$ un homomorphisme local, $M$ un $B$-module de type fini.
Soit $(f_i)_{1\leq i\leq n}$ une suite d'éléments de l'idéal maximal de $B$,
et pour tout $i$, soit $g_i$ l'image de $f_i$ dans $B\otimes_A k$. Les
conditions suivantes sont équivalentes~:
\begin{enumerate}
  \item[\rm(a)] La suite $(f_i)$ est $M$-régulière et les quotients
  \[
    M_i=M/\left(\sum_{j=1}^i f_jM\right)
  \]
  sont des $A$-modules plats pour $1\leq i\leq n$.
  \item[\rm(b)] La suite $(f_i)$ est $M$-régulière et le $A$-module $M_n$ est
  plat.
  \item[\rm(c)] Le $A$-module $M$ est plat et la suite
  $(g_i)_{1\leq i\leq n}$ est $(M\otimes_A k)$-régulière.
  \item[\rm(d)] Le $A$-module $M$ est plat, et pour tout homomorphisme
  d'anneaux $\rho:A\to A'$, si on pose $B'=B\otimes_A A'$,
  $M'=M\otimes_A A'$, $f'_i=f_i\otimes1$ ($1\leq i\leq n$), la suite
  $(f'_i)$ est $M'$-régulière.
\end{enumerate}
\end{proposition}

\begin{proof}
Il est clair que a) entraîne b)~; pour montrer que b) entraîne d), il suffit,
en vertu de (15.1.15), de montrer que b) entraîne que les $M_i$ sont des
$A$-modules plats pour $0\leq i\leq n-1$ (en posant $M_0=M$). Comme
$M_{i+1}=M_i/f_{i+1}M_i$, on est ramené, par récurrence descendante, à prouver
que si $h$ est un élément de l'idéal maximal de $B$ tel que l'homothétie de
rapport $h$ dans $M$ soit injective et que $M/hM$ soit un $A$-module plat,
alors $M$ est un $A$-module plat~; mais cela n'est autre que
($0_{\mathrm{III}}$, 10.2.7). Il est clair que c) est le cas particulier de
d) obtenu en prenant $A'=k$. Reste donc à prouver que c) entraîne a). Comme
$M$ est un $A$-module plat et un $B$-module de type fini, il résulte de c) et
de ($0_{\mathrm{III}}$, 10.2.4) que l'homothétie de rapport $f_1$ dans $M$
est injective et que son conoyau $M/f_1M=M_1$ est un $A$-module plat.
Raisonnons par récurrence sur $i$ et supposons démontré que $M_i$ est un
$A$-module plat~; comme l'homothétie de rapport $g_{i+1}$ dans
$M_i\otimes_A k$ est injective par hypothèse, il résulte encore de
($0_{\mathrm{III}}$, 10.2.4) que l'homothétie de rapport $f_{i+1}$ dans
$M_i$ est injective et que $M_i/f_{i+1}M_i=M_{i+1}$ est un $A$-module plat,
ce qui achève la démonstration. On notera que la démonstration du fait que c)
entraîne a) n'utilise pas le fait que les $f_i$ appartiennent à l'idéal
maximal de $B$.
\end{proof}

\begin{corollary}[15.1.17]
\label{0.15.1.17-fr}
Les hypothèses sur $A$ et $B$ étant celles de (15.1.16), soient $M$ un
$B$-module de type fini, plat sur $A$, $R$ un sous-$B$-module de $M$ tel que
$N=M/R$ soit un $A$-module plat. Soit $(f_i)_{1\leq i\leq n}$ une suite
d'éléments de l'idéal maximal de $B$ ayant les propriétés suivantes~:
\begin{enumerate}
  \item[(i)] $f_iM\subset R$ pour $1\leq i\leq n$.
  \item[(ii)] Si $g_i$ est l'image de $f_i$ dans $B\otimes_A k$, la suite
  $(g_i)_{1\leq i\leq n}$ est $(M\otimes_A k)$-régulière.
  \item[(iii)] La somme des sous-$(B\otimes_A k)$-modules
  $g_i(M\otimes_A k)$ de $M\otimes_A k$ est l'image canonique de
  $R\otimes_A k$ dans $M\otimes_A k$ (égale au noyau de l'homomorphisme
  canonique $M\otimes_A k\to N\otimes_A k$).
\end{enumerate}
Alors la suite $(f_i)$ est $M$-régulière et l'on a
$R=\sum_{i=1}^n f_iM$.
\end{corollary}

\begin{proof}
\oldpage[0\textsubscript{IV-1}]{19}
Comme $N$ est un $A$-module plat, il résulte de l'exactitude de la suite
$0\to R\to M\to N\to0$ que la suite
\[
  0\longrightarrow R\otimes_A k\longrightarrow M\otimes_A k
  \longrightarrow N\otimes_A k\longrightarrow0
\]
est exacte ($0_{\mathrm I}$, 6.1.2)~; on a donc
$R\otimes_A k=\sum_i g_i(M\otimes_A k)$. Comme $M$ est un $B$-module de type
fini et que $B$ est noethérien, $R$ est aussi un $B$-module de type fini,
d'où $R=\sum_i f_iM$ en vertu du lemme de Nakayama, les $f_i$ appartenant à
l'idéal maximal de $B$.
\end{proof}

\begin{lemma}[15.1.18]
\label{0.15.1.18-fr}
Soient $A$ un anneau, $M$ un $A$-module muni d'une filtration séparée et
exhaustive $(M_k)_{k\geq0}$, et soit $f\in A$ un élément
$\gr_\bullet(M)$-régulier. Alors~:
\begin{enumerate}
  \item[(i)] $f$ est $M$-régulier.
  \item[(ii)] Pour tout $k$, on a $M_k\cap fM=fM_k$, et l'image canonique
  $M'_k$ de $M_k$ dans $M'=M/fM$ est donc canoniquement isomorphe à
  $M_k/fM_k$.
  \item[(iii)] La suite
  \[
    0\longrightarrow M_{k+1}/fM_{k+1}\longrightarrow M_k/fM_k
    \longrightarrow\gr_k(M)/f\cdot\gr_k(M)\longrightarrow0
  \]
  est exacte, et si l'on munit $M'$ de la filtration quotient $(M'_k)$,
  $\gr_\bullet(M')$ est donc canoniquement isomorphe à
  $\gr_\bullet(M)/f\cdot\gr_\bullet(M)$.
\end{enumerate}
\end{lemma}

\begin{proof}
L'assertion (i) résulte de l'hypothèse sur $f$ et de Bourbaki,
\emph{Alg.~comm.}, chap.~III, \S~2, no~8, cor.~1 du th.~1. Comme
$\gr_\bullet(M/M_k)$ est facteur direct de $\gr_\bullet(M)$, $f$ est
$\gr_\bullet(M/M_k)$-régulier pour tout $k$, donc $(M/M_k)$-régulier par
(i), ce qui établit (ii). Pour prouver (iii), il suffit de voir que
l'application $M_{k+1}/fM_{k+1}\to M_k/fM_k$ est injective~: or, si
$x\in M_{k+1}$ est tel que $x\in fM_k$, il résulte de (ii) que
$x\in fM_{k+1}$, ce qui achève de prouver le lemme.
\end{proof}

\begin{proposition}[15.1.19]
\label{0.15.1.19-fr}
Soient $A$ un anneau, $M$ un $A$-module muni d'une filtration
$(M_k)_{k\geq0}$ telle que $M_0=M$, $(f_i)_{1\leq i\leq n}$ une suite
$\gr_\bullet(M)$-régulière d'éléments de $A$. Supposons en outre que pour
$0\leq i\leq n-1$, la filtration quotient de $(M_k)$ sur
\[
  M/\left(\sum_{j\leq i}f_jM\right)
\]
soit séparée. Alors la suite $(f_i)$ est $M$-régulière, et si l'on pose
\[
  M'=M/\left(\sum_{i\leq n}f_iM\right)
\]
et que l'on munit $M'$ de la filtration quotient de $(M_k)$,
$\gr_\bullet(M')$ est canoniquement isomorphe à
\[
  \gr_\bullet(M)/
  \left(\sum_{i\leq n}f_i\gr_\bullet(M)\right).
\]
\end{proposition}

\begin{proof}
Pour $n=1$, la proposition résulte du lemme (15.1.18). Raisonnons par
récurrence sur $n$, en posant
\[
  N=M/\left(\sum_{i\leq n-1}f_iM\right)
\]
et en désignant par $(N_k)$ la filtration quotient de $(M_k)$ sur $N$, qui
est séparée par hypothèse~; en outre $\gr_\bullet(N)$ est isomorphe à
\[
  \gr_\bullet(M)/
  \left(\sum_{i\leq n-1}f_i\gr_\bullet(M)\right).
\]
Par hypothèse $f_n$ est donc $\gr_\bullet(N)$-régulier, et il résulte donc du
lemme (15.1.18) appliqué à $N$ que $f_n$ est $N$-régulier et
$\gr_\bullet(M')$ isomorphe à
$\gr_\bullet(N)/f_n\gr_\bullet(N)$, ce qui démontre la proposition.
\end{proof}

La proposition (15.1.19) s'appliquera en particulier pour les filtrations
vérifiant l'une ou l'autre hypothèse suivante~:
\begin{enumerate}
  \item[$1^\circ$] La filtration $(M_k)$ est finie et séparée (puisque cela
  entraîne $M_k=0$ pour $k$ assez grand)~;
  \item[$2^\circ$] $A$ est un anneau noethérien, $\mathfrak{J}$ un idéal
  contenu dans le radical de $A$, $M$ un $A$-module de type fini et $(M_k)$
  la filtration $\mathfrak{J}$-préadique ($0_{\mathrm I}$, 7.3.5).
\end{enumerate}

\begin{corollary}[15.1.20]
\label{0.15.1.20-fr}
\oldpage[0\textsubscript{IV-1}]{20}
Soient $A$ un anneau, $M$ un $A$-module, $(f_i)_{1\leq i\leq n}$ une suite
$M$-régulière d'éléments de $A$. Alors, quels que soient les entiers
$\nu_i\geq1$, la suite $(f_i^{\nu_i})$ est $M$-régulière et
$M/(\sum_{i\leq n}f_i^{\nu_i}M)$ a une suite de composition dont les
quotients sont isomorphes à $M/(\sum_{i\leq n}f_iM)$.
\end{corollary}

\begin{proof}
Raisonnons par récurrence sur $n$~; pour $n=1$ les assertions sont immédiates
puisque $f^jM/f^{j+1}M$ est isomorphe à $M/fM$, l'homothétie de rapport $f$
dans $M$ étant injective par hypothèse. Posons alors
\[
  N=M/\left(\sum_{i\leq n-1}f_i^{\nu_i}M\right)~;
\]
par l'hypothèse de récurrence, $N$ admet une filtration finie $(N_k)$ pour
laquelle les $\gr_k(N)$ sont isomorphes à
$M/(\sum_{i\leq n-1}f_iM)$~; par hypothèse $f_n$ est donc
$\gr_\bullet(N)$-régulier~; il résulte par suite de (15.1.19) que $f_n$ est
$N$-régulier. Appliquant le cas $n=1$ démontré au début, on en conclut que
$f_n^{\nu_n}$ est $N$-régulier et que $N/f_n^{\nu_n}N$ admet une suite de
composition dont les quotients sont isomorphes à $N/f_nN$~; mais pour la
filtration de $N/f_nN$ quotient de la filtration $(N_k)$,
$\gr_k(N/f_nN)$ est isomorphe à $\gr_k(N)/f_n\gr_k(N)$ par (15.1.19) pour
tout $k$, donc à $M/(\sum_{i\leq n}f_iM)$, ce qui prouve le corollaire.
\end{proof}

\begin{proposition}[15.1.21]
\label{0.15.1.21-fr}
Soient $A$, $B$ deux anneaux locaux noethériens, $\varphi:A\to B$ un
homomorphisme local, $M$ un $B$-module de type fini, $(f_i)_{1\leq i\leq n}$
une suite $A$-régulière formée d'éléments de l'idéal maximal de $A$. Les
propriétés suivantes sont équivalentes~:
\begin{enumerate}
  \item[a)] $M$ est un $A$-module plat.
  \item[b)] Si $\mathfrak{J}$ est l'idéal $\sum_{i\leq n}f_iA$,
  $M/\mathfrak{J}M$ est un $A/\mathfrak{J}$-module plat et la suite $(f_i)$
  est $M$-régulière.
\end{enumerate}
\end{proposition}

\begin{proof}
Le fait que a) entraîne b) résulte de (15.1.13) et de
($0_{\mathrm{III}}$, 10.2.2). Inversement, si la suite $(f_i)$ est
$M$-régulière, elle est $M$-quasi-régulière (15.1.9), donc l'homomorphisme
(15.1.1.1) est bijectif. Mais comme la suite $(f_i)_{1\leq i\leq n}$ est aussi
supposée $A$-régulière, elle est $A$-quasi-régulière (15.1.9), et
l'homomorphisme canonique (15.1.1.1) correspondant
$(A/\mathfrak{J})[T_1,\ldots,T_n]\to\gr_\bullet(A)$ est bijectif. On en
conclut que l'homomorphisme canonique ($0_{\mathrm{III}}$, 10.1.1.2)
\[
  \gr_0(M)\otimes_{A/\mathfrak{J}}\gr_\bullet(A)
  \longrightarrow\gr_\bullet(M)
\]
est bijectif. La conclusion résulte alors de ($0_{\mathrm{III}}$, 10.2.2),
les $f_i$ étant dans l'idéal maximal de $A$ et l'homomorphisme $\varphi$ étant
local.
\end{proof}

\subsection{Suites \texorpdfstring{$\mathscr{F}$}{F}-régulières}
\label{subsection:0.15.2-fr}

\begin{env}[15.2.1]
\label{0.15.2.1-fr}
Soient $X$ un espace annelé, $\mathscr{F}$ un $\mathcal{O}_X$-Module,
$f_i$ ($1\leq i\leq n$) une suite de sections de $\mathcal{O}_X$ au-dessus
de $X$, $\mathscr{J}$ l'Idéal de $\mathcal{O}_X$ engendré par les $f_i$
(image de $\mathcal{O}_X^n$ par l'homomorphisme
$\mathcal{O}_X^n\to\mathcal{O}_X$ défini canoniquement par les $f_i$
($0_{\mathrm I}$, 5.1.1)). Les sous-$\mathcal{O}_X$-Modules
$\mathscr{J}^k\mathscr{F}$ définissent sur $\mathscr{F}$ une filtration, et
on pose encore
$\gr_k(\mathscr{F})=\mathscr{J}^k\mathscr{F}/\mathscr{J}^{k+1}\mathscr{F}$~;
si $S_k$ est l'ensemble des polynômes homogènes de degré total $k$ dans
l'anneau $\mathbb{Z}[T_1,\ldots,T_n]$, $S_k$ peut être considéré comme un
faisceau simple sur $X$, et on définit comme dans (15.1.1) un homomorphisme
canonique
\[
\label{0.15.2.1.1-fr}
  \varphi_k:\gr_0(\mathscr{F})\otimes_{\mathbb Z}S_k
  \longrightarrow\gr_k(\mathscr{F}).
\tag{15.2.1.1}
\]
\oldpage[0\textsubscript{IV-1}]{21}
de la façon suivante~: pour tout ouvert $U$ de $X$ (ou d'une base de la
topologie de $X$), on considère l'homomorphisme (15.1.1.1)
\[
  \varphi_{k,U}:\gr_0(\Gamma(U,\mathscr{F}))\otimes_{\mathbb Z}S_k
  \longrightarrow\gr_k(\Gamma(U,\mathscr{F}))
\]
défini par les $n$ sections $f_i|U\in\Gamma(U,\mathcal{O}_X)$, le
$\Gamma(U,\mathcal{O}_X)$-module $\Gamma(U,\mathscr{F})$ étant filtré par les
sous-modules $(\Gamma(U,\mathscr{J}))^k\Gamma(U,\mathscr{F})$. Le
$\mathcal{O}_X$-Module $\gr_k(\mathscr{F})$ étant associé au préfaisceau
$U\mapsto\gr_k(\Gamma(U,\mathscr{F}))$ pour tout $k$, les $\varphi_{k,U}$
définissent un homomorphisme de $\mathcal{O}_X$-Modules (15.2.1.1). En vertu
des propriétés d'exactitude des limites inductives de modules et de leur
commutation aux produits tensoriels, il résulte de ce qui précède que pour
tout $x\in X$, l'homomorphisme $(\varphi_k)_x$ des fibres des deux membres de
(15.2.1.1) au point $x$ s'identifie à l'homomorphisme de (15.1.1)
\[
\label{0.15.2.1.2-fr}
  \gr_0(\mathscr{F}_x)\otimes_{\mathbb Z}S_k
  \longrightarrow\gr_k(\mathscr{F}_x)
\tag{15.2.1.2}
\]
défini par la suite $((f_i)_x)_{1\leq i\leq n}$ d'éléments de
$\mathcal{O}_x$.
\end{env}

\begin{env}[15.2.2]
\label{0.15.2.2-fr}
Les notations étant celles de (15.2.1), nous dirons que la suite
$(f_i)_{1\leq i\leq n}$ est $\mathscr{F}$-régulière si, en désignant par
$f_j\mathscr{F}$ le sous-$\mathcal{O}_X$-Module de $\mathscr{F}$ image de
l'endomorphisme de $\mathscr{F}$ défini par $f_j$, l'endomorphisme de
$\mathscr{F}/(\sum_{j=1}^{i-1}f_j\mathscr{F})$ défini par $f_i$ est injectif
pour $1\leq i\leq n$~; nous dirons que la suite $(f_i)$ est
$\mathscr{F}$-quasi-régulière si les homomorphismes (15.2.1.1) sont injectifs.
Il résulte de ces définitions et des remarques de (15.2.1) que pour que la
suite $(f_i)$ soit $\mathscr{F}$-régulière (resp.
$\mathscr{F}$-quasi-régulière), il faut et il suffit que, pour tout $x\in X$,
la suite $((f_i)_x)_{1\leq i\leq n}$ soit $\mathscr{F}_x$-régulière (resp.
$\mathscr{F}_x$-quasi-régulière). Toute suite $\mathscr{F}$-régulière est donc
$\mathscr{F}$-quasi-régulière (15.1.9)~; la réciproque est vraie si
$\mathscr{F}$ est un $\mathcal{O}_X$-Module de type fini, si les
$\mathcal{O}_x$ sont des anneaux noethériens pour tout $x\in X$ et si en outre
les $(f_i)_x$ appartiennent au radical de $\mathcal{O}_x$ pour tout $x\in X$
(15.1.11).

Si $X=\operatorname{Spec}(A)$ est un schéma affine, et
$\mathscr{F}=\widetilde{M}$, il
résulte de (I, 1.3) que pour que la suite $(f_i)$ soit
$\mathscr{F}$-régulière (resp. $\mathscr{F}$-quasi-régulière), il faut et il
suffit qu'elle soit $M$-régulière (resp. $M$-quasi-régulière).

Lorsque $\mathscr{F}=\mathcal{O}_X$, on supprimera la mention de
$\mathscr{F}$ dans les définitions précédentes.
\end{env}

\begin{remark}[15.2.3]
\label{0.15.2.3-fr}
Lorsque $X=\operatorname{Spec}(A)$ et
$\mathfrak{J}=\sum_{i=1}^n f_iA$, la suite $(f_j)$ est
quasi-régulière pour tout $A$-module $M$ dans tout ouvert $U$ ne rencontrant
pas $V(\mathfrak{J})$, car en tout point $x\in U$, on a
$\mathfrak{J}_x=\mathcal{O}_x$ et les deux membres de (15.2.1.2) sont nuls. La
notion de suite $\widetilde{M}$-régulière n'a donc d'intérêt qu'au voisinage
des points de $V(\mathfrak{J})$.
\end{remark}

\begin{proposition}[15.2.4]
\label{0.15.2.4-fr}
Soient $X$ un espace annelé, $\mathscr{F}$ un $\mathcal{O}_X$-Module cohérent,
$x$ un point de $X$, $(f_i)_{1\leq i\leq n}$ une suite d'éléments de
$\Gamma(X,\mathcal{O}_X)$. Si la suite $((f_i)_x)_{1\leq i\leq n}$ est
$\mathscr{F}_x$-régulière, il existe un voisinage ouvert $U$ de $x$ tel que la
suite $(f_i|U)_{1\leq i\leq n}$ soit $(\mathscr{F}|U)$-régulière.
\end{proposition}

\begin{proof}
En effet, pour $1\leq i\leq n$, le $\mathcal{O}_X$-Module
$\mathscr{F}/(\sum_{j=1}^{i-1}f_j\mathscr{F})$ est cohérent
($0_{\mathrm I}$, 5.3.4), donc il en est de même du noyau de
l'endomorphisme de ce $\mathcal{O}_X$-Module défini par $f_i$
($0_{\mathrm I}$, 5.3.4)~; comme la fibre au point $x$ de ce noyau est nulle,
il en est de même de sa restriction à un voisinage de $x$
($0_{\mathrm I}$, 5.2.2).
\end{proof}

\begin{proposition}[15.2.5]
\label{0.15.2.5-fr}
\oldpage[0\textsubscript{IV-1}]{22}
Soit $u=(\psi,\theta):X\to Y$ un morphisme plat d'espaces annelés
($0_{\mathrm I}$, 6.7.1)~; soient $\mathscr{F}$ un
$\mathcal{O}_Y$-Module, $\mathscr{G}=u^*(\mathscr{F})$,
$(f_i)_{1\leq i\leq n}$ une suite d'éléments de
$\Gamma(Y,\mathcal{O}_Y)$, $(g_i)_{1\leq i\leq n}$ la suite de leurs images
$g_i=\theta(f_i)\in\Gamma(Y,u_*(\mathcal{O}_X))=\Gamma(X,\mathcal{O}_X)$. Si
la suite $(f_i)$ est $\mathscr{F}$-régulière, la suite $(g_i)$ est
$\mathscr{G}$-régulière. La réciproque est vraie si $u$ est un morphisme
fidèlement plat ($0_{\mathrm I}$, 6.7.8).
\end{proposition}

\begin{proof}
Soient $x$ un point de $X$, $y=u(x)=\psi(x)$~; on a
\[
  \mathscr{G}_x=\psi_x^*(\mathscr{F}_y)
  \otimes_{\psi_x^*(\mathcal{O}_y)}\mathcal{O}_x,
  \qquad\text{et}\qquad
  (g_i)_x=\theta_x^\#(\psi_x^*((f_i)_y))~;
\]
comme par hypothèse
$\theta_x^\#:\psi_x^*(\mathcal{O}_y)\to\mathcal{O}_x$ fait de
$\mathcal{O}_x$ un $\psi_x^*(\mathcal{O}_y)$-module plat, et que le foncteur
$\psi^*$ est exact, la première assertion résulte de l'assertion correspondante
de (15.1.14)~; il en est de même de la seconde, compte tenu de ce que
l'hypothèse que $u$ est fidèlement plat signifie que $\psi$ est surjectif et
que $\theta_x^\#$ fait de $\mathcal{O}_x$ un
$\psi_x^*(\mathcal{O}_y)$-module fidèlement plat.
\end{proof}
